\chapter{Ordinamento dei risulati}
\section{ORDER BY}
\subsection{Cosa è?}
L'ORDER BY è una clausola opzionale che si inserisce al termina di una query SELECT per ordinare le righe del risultato in baso ai valori di una o più colonne
\subsection{Sintassi}
\begin{lstlisting}[language=sql_generic]
SELECT colonna1, colonna2, ...
FROM NomeTabella
[WHERE condizioni_di_filtro]
ORDER BY colonnaA [ASC | DESC], colonnaB [ASC | DESC];
\end{lstlisting}
\subsubsection{Ordinamento}
\begin{itemize}
    \item \textbf{Ascendente (ASC)} è default e può essere omesso
    \item \textbf{Discendente (DESC)}
\end{itemize}
\subsection{Esempio}
\begin{lstlisting}[language=sql_generic]
SELECT 
    UCASE(I.Nome) AS Nome_Impiegato,
    ((I.Stipendio + I.Premio) * 13) AS Guadagno_Annuo_Stimato,
    D.Citta AS Sede_Dipartimento
FROM Impiegato AS I
JOIN Dip AS D 
  ON I.Dipart = D.Nome
WHERE UCASE(I.Mansione) = 'INGEGNERE'
ORDER BY Guadagno_Annuo_Stimato DESC, Nome_Impiegato ASC;
\end{lstlisting}

\section{Operatori di aggregazione}
\subsection{Quali sono?}
\begin{itemize}
    \item \textbf{COUNT:} conta numero di tuple di tabella
    \item \textbf{SUM:} somma i valori o espressioni di attributi
    \item \textbf{MAX:} valore massimo di un attributo di tabella
    \item \textbf{MIN:} valore minimo di un attributo di tabella
    \item \textbf{AVG:} valore medio di un attributo di tabella
    \item \textbf{GROUP BY:} raggruppamento delle tuple in sottogrupi gestiti come estensione delle normali interrogazioni
\end{itemize}
\subsection{Come vengono applicati?}
\begin{enumerate}
    \item Si esegue l'interrogazione sulla base di clausole FROM e WHERE
    \item Si applica l'operatore aggregato alla tabella risultato dell'interrogazione
\end{enumerate}

\section{COUNT}
\subsection{Cosa fa?}
Restituisce il numero di valori degli attributi in ListaAttr
\subsection{Sintassi}
\begin{lstlisting}[language=sql_generic]
COUNT (< * | [DISTINCT | ALL] ListaAttr >)
\end{lstlisting}
\begin{itemize}
    \item \textbf{DISTINCT:} numero di valori distinti diversi da NULL di ListaAttr
    \item \textbf{ALL:} numero di valori diversi da NULL di ListaAttr
    \item \textbf{COUNT (*):} conta numero di righe di una tabella
\end{itemize}
Se non è specificato nulla, ALL è l'opzione di default

\subsection{Esempio}
\begin{lstlisting}[language=sql_generic]
SELECT 
    D.Nome AS Nome_Dipartimento,
    D.Citta AS Sede,
    COUNT(I.Matricola) AS Numero_Impiegati
FROM Dip AS D
JOIN Impiegato AS I 
  ON D.Nome = I.Dipart
WHERE I.Stipendio > 1500
GROUP BY D.Nome, D.Citta
HAVING COUNT(I.Matricola) > 2
ORDER BY Numero_Impiegati DESC;
\end{lstlisting}

\section{SUM}
\subsection{Cosa fa?}
Restituisce la somma dei valori di AttrExpr relativamente alle tuple di tabella specificate nella query, opera su attributi numerici
\subsection{Sintassi}
\begin{lstlisting}[language=sql_generic]
SELECT SUM(nome_colonna) AS nome_alias
FROM NomeTabella
[WHERE condizioni_di_filtro];
\end{lstlisting}
\subsection{Esempio}
\begin{lstlisting}[language=sql_generic]
SELECT SUM(Stipendio) AS Spesa_Totale_Stipendi
FROM Impiegato;
\end{lstlisting}

\section{MIN/MAX}
\subsection{Cosa fa?}
Restituisce il massimo/minimo dei valori di AttrExpr relativamente alle tuple di tabella specificate nella query, opera su attributi  ordinabili
\subsection{Sintassi}
\begin{lstlisting}[language=sql_generic]
SELECT 
    MIN(nome_colonna) AS minimo_alias,
    MAX(nome_colonna) AS massimo_alias
FROM NomeTabella
[WHERE condizioni_di_filtro];
\end{lstlisting}
\subsection{Esempio}
\begin{lstlisting}[language=sql_generic]
SELECT 
    MIN(Stipendio) AS Stipendio_Minimo,
    MAX(Stipendio) AS Stipendio_Massimo
FROM Impiegato;
\end{lstlisting}

\section{AVG}
\subsection{Cosa fa?}
Restituisce la media dei valori di AttrExpr relativamente alle tuple di tabella specificate nella query, opera su valori numerici
\subsection{Sintassi}
\begin{lstlisting}[language=sql_generic]
SELECT AVG(nome_colonna) AS nome_alias
FROM NomeTabella
[WHERE condizioni_di_filtro];
\end{lstlisting}
\subsection{Esempio}
\begin{lstlisting}[language=sql_generic]
SELECT AVG(Stipendio) AS Stipendio_Medio_Aziendale
FROM Impiegato;
\end{lstlisting}

\section{Target List}
SQL impedisce di includere in una stessa target list funzioni aggregate e espressioni al livello di riga

\begin{lstlisting}[language=sql_generic]
SELECT Cognome, MAX(Stipendio) -- Questa cosa non si puo fare
FROM Impiegato
WHERE Dipart = "Amministrazione"
\end{lstlisting}

È necessario ricorrere alle subquery, che si vedranno più avanti

\section{Operatore di aggregamento GROUP BY}
\subsection{Cosa è?}
La GROUP BY è una clausola dell'istruzione SELECT utilizzata per raggruppare le righe di una tabella che condividono gli stessi valori in una o più colonne specificate.
\subsection{Sintassi}
\begin{lstlisting}[language=sql_generic]
SELECT colonna_di_raggruppamento, FUNZIONE_AGGREGATA(altra_colonna)
FROM NomeTabella
[WHERE filtri_sulle_singole_righe]
GROUP BY colonna_di_raggruppamento
[HAVING filtri_sui_gruppi_creati]
[ORDER BY criterio_di_ordinamento];
\end{lstlisting}
\subsection{Esempio}
\begin{lstlisting}[language=sql_generic]
SELECT Dipart, AVG(Stipendio) AS Stipendio_Medio
FROM Impiegato
GROUP BY Dipart;
\end{lstlisting}

\section{Condizioni WHERE e HAVING}
Specificano entrambi condizioni di selezione:
\begin{itemize}
    \item \textbf{WHERE:} Predicati che operano su tuple individuabili (seleziona di tuple prima di GROUP BY)
    \item \textbf{HAVING:} Predicati che operano su gruppi di tuple creati da GROUP BY (selezioni raggruppamenti)
\end{itemize}

\subsection{Sintassi WHERE}
\begin{lstlisting}[language=sql_generic]
WHERE attributo OPERATORE valore_o_costante
  [ AND | OR | NOT altra_condizione ];
\end{lstlisting}

\subsection{Esempio di uso di WHERE}
\begin{lstlisting}[language=sql_generic]
SELECT Cognome, Stipendio
FROM Impiegato
WHERE Dipart = 'Amministrazione' AND Stipendio > 1500;
\end{lstlisting}

\subsection{Sintassi HAVING}
\begin{lstlisting}[language=sql_generic]
GROUP BY colonna_raggruppata
HAVING FUNZIONE_AGGREGATA(colonna) OPERATORE valore;
\end{lstlisting}

\subsection{Esempio di uso di HAVING}
\begin{lstlisting}[language=sql_generic]
SELECT Dipart, AVG(Stipendio) AS AvgStipendio
FROM Impiegato
GROUP BY Dipart
HAVING COUNT(*) >= 2; -- Mantiene solo i dipartimenti con almeno 2 impiegati
\end{lstlisting}